\chapter{June}

\section[June 1st]{June 1\textsuperscript{st}}
The class started with Professor Cyntia asking everyone individually how they were doing. We could not repeat ourselves. I find this idea of Professor Cyntia not allowing us to simply say ``Fine'' interesting; it strikes me as strange, and it feels a bit uncomfortable to have to say that I am tired, happy, or anything else.

I feel like a really lazy person, or like I am complaining, every time I say that I am ``tired'' or ``sleepy''. This is a bit morally charged; it sounds like I am complaining or ``victimizing'' myself. When I said I was sleepy, the professor stated that she ``hopes not to put me to bed during class''. Of course, it probably feels uncomfortable for her to know that she has a ``sleepy'' student.

This activity really makes me value how well society is organized so that we all agree to not actually express how we really are. A simple ``I'm fine'' or ``I'm good'' shows politeness without being invasive. When a person does not accept answers such as ``I'm fine'' or ``I feel normal'' and demands a more meaningful answer, it simply feels rude, as if, somehow, I need to answer to her regarding how I am doing.

The book for today was ``Be You''; the class started with a little bit of vocabulary. It is about how different adjectives relate to a child, highlighting many positive characteristics such as being loved, effervescent, connected, brave, okay with being alone, and persistent.

One point of discussion was ``why children are so curious, and when we grow up we lose this curiosity about asking why.'' My classmates suggested that routine and work make us tired and distracted, removing this curiosity from us. I disagree with that. It seems that, in reality, we merely get all of our questions answered as we grow up. We learn where the rain comes from, how fire works, what a lemon tastes like, what happens when we stub our toes, and so on. It turns out that the answers to our questions were hardly that exciting. As we live, we gather all the knowledge we need for our daily lives; once we learn the basics, there are only complex questions left that demand deeper study and research, and the work of learning those complex topics just isn't worth the effort sometimes.

It is easy to be amazed by simple things before we learn about them. As a child, I loved playing with water, until I started washing the dishes every day. Now, water is quite boring. The same goes for fire: lighting a match and seeing the flames was fun, but it just got old after a while. There comes a time when we eventually find the answers to all our questions, and they run out. If we force ourselves to learn new things that we don't really want to learn, like instruments, sports, and so on, we start to accumulate unfinished courses and books that will never be completed, because the questions they answer are no longer relevant to us.

Another point of discussion was about being ``ourselves,'' reflecting upon how people sometimes tend to alter their behavior in a group of friends just to fit in. I would say that there is quite a bit of discussion to be had regarding the term ``self''. It can be surprisingly difficult to determine who we really are, as it is impossible to isolate ourselves from external influences. We are shaped by our upbringing at home, by our teachers, by the news, by social media, our internal moral code, what we study, and many other factors. It reaches a point where it becomes remarkably difficult to determine a static ``self'' to be preserved, and the friends we have are just another variable in that mix.

I understand the point being made is that we shouldn't force ourselves to do things we don't want to do, especially not purely to impress others. We should be critical about our social relationships. Sometimes it is beneficial to change; our friends can teach us to be better individuals, provide examples of integrity, responsibility, and intelligence, show different perspectives, and offer constructive criticism. The opposite is also true: negative friendships can provide bad examples and make our lives considerably worse.

The professor says that she uses this book in class to build ``repertoire''. She asks, ``how can we be good writers if we cannot be good speakers and talk about ourselves?'' I can see the point, although it feels like stretching the idea a bit. The content of the book just feels overly simplistic. It is really difficult to improve my ``repertoire'' through discussions on these books; in fact, it feels like trying to get blood out of a stone. We could be reading abstracts and academic summaries on psychology, sociology, and philosophy regarding the themes of these books. In that way, we would learn better ways to express ourselves academically while discussing the topics, perhaps with a more critical eye, going beyond simple messages made for children. A more compelling argument might be that the class is composed of teachers and future teachers, so they need to learn to read books and talk about them with children. Even then, it doesn't seem like this is the best class for that, as it was already covered in the ``English for Kids'' class.

Sure, such books generated a lot of material for this reflective portfolio; probably the majority of this portfolio is about topics not related to ``Academic Writing'' at all. Maybe this is indicative of the tangents that end up happening quite a bit during class.

The next topic was about abstracts themselves. When I wrote my first little article in high school, it felt like just an obligation. Now in university, I see their utility; I view them mainly as a ``screening device''. When I need to find articles to base my research on, it helps a lot to follow a specific process: read the title; if it is interesting, read the abstract; if that is also good, read the conclusion or results section. Those three elements are the main determinants I use to decide if a paper is worth my time. If the title isn't interesting, I drop it, and the same goes for the abstract and, finally, the conclusion.

Bringing it to the philosophical term of ``abstraction'', in Aristotelian philosophy, abstraction is treated as a special subtraction. It happens when, through experience, we can subtract the essence of something. For example, after seeing many chairs, we can subtract the ``accidents'', like the materials and the position, and realize that, besides the differences, we can understand that they are chairs.

The abstract of an article demands a similar approach: subtracting most of the content of the article while keeping only the ``essence'': what it truly is and what makes it unique.


\section[June 8th]{June 8\textsuperscript{th}}

Today's book is called ``Perfect''. It is about a pencil that makes drawings, irritating the eraser that runs after the pencil erasing the squiggles and drawings made by the pencil. The pencil ended up making other pencils and dominated everything with drawings, while the eraser got thinner and thinner. Eventually, the eraser started to do ``drawings'' by erasing the page that was fully covered by the pencil. After some time, the eraser found a blank part and got sad because it couldn't draw anymore. At the end, the eraser called the pencil again so that it could draw everywhere again, and the eraser could make its drawings by erasing the covered page. The message of the book appears to be about optimism and against being persnickety (someone who is fussy about small details).

After the reading, we reflected upon the concept of ``perfect''. My colleagues mentioned things such as ``lasagnas'', ``things that are not common'', or even ``flawless''. The idea is that perfection is subjective; every person has their own definition of perfect, and such definitions don't match, making it necessary to make compromises. The professor called this discussion philosophical.

It seems that there is actually a more precise definition of perfect. A truly philosophical one: it is completeness in actuality. ``A thing is perfect when it possesses all the properties or attributes proper to its nature and completely lacks unrealized potential''.

After these discussions, we reviewed the 5 moves, which I have already discussed at length in this portfolio.

Following the 5 moves review, we returned to check the abstracts that the class had collected.


\begin{quote}
    \textbf{From English as a Lingua Franca to Englishes as Linguae Harmonicae: A Critical Framework for Ethical and Relational Language Education}

    This article introduces Englishes as Linguae Harmonicae (ESLH), a new theoretical and pedagogical framework that reimagines English language education through the lenses of relational ethics, emotional resonance, and decolonial care. Building on and moving beyond English as a Lingua Franca (ELF), World Englishes (WEs), and Global Englishes (GEs), ESLH foregrounds harmony not as sameness, but as co-resonance across difference. Rooted in pluriversal and postcolonial thought, ESLH emphasizes attunement, dignity, and healing in response to the postcolonial, emotional, and epistemic harms of English. Drawing on metaphors of music and restorative practice, the article proposes a praxiology of English teaching centered on relational presence, dialogical listening, and ethical co-creation. ESLH invites educators to imagine classrooms not as standardized spaces, but as soundscapes of belonging. In reframing Englishes as sites of justice, mutuality, and shared humanity, this framework contributes to a future of language education grounded in critical love and deeper understanding.
    
    (149 words)
\end{quote}

This abstract REALLY looks AI-generated to me. When analyzing AI-generated texts, we usually check for 3 types of constructions:
\begin{enumerate}
    \item Arbitrary lists of three elements.
    \item Excessive constructions of formats similar to ``It is not only X, but Y''.
    \item Usage of exaggerated and excessively dramatic adjectives for things that are really not that big of a deal.
\end{enumerate}

I can find plenty of all three points in the text. Here are the examples:

\begin{enumerate}
    \item \textbf{Arbitrary lists of three elements:}
    \begin{itemize}
        \item ``through the lenses of relational ethics (1), emotional resonance (2), and decolonial care (3);''
        \item ``Building on and moving beyond English as a Lingua Franca (ELF) (1), World Englishes (WEs) (2), and Global Englishes (GEs) (3);''
        \item ``ESLH emphasizes attunement (1), dignity (2), and healing (3);''
        \item ``in response to the postcolonial (1), emotional (2), and epistemic harms (3) of English;''
        \item ``the article proposes a praxiology of English teaching centered on relational presence (1), dialogical listening (2), and ethical co-creation (3);''
        \item ``In reframing Englishes as sites of justice, mutuality, and shared humanity.''
    \end{itemize}
    
    \item \textbf{Excessive constructions of formats similar to ``It is not only X, but Y'':}
    \begin{itemize}
        \item ``ESLH foregrounds harmony not as sameness, but as co-resonance across difference''
        \item ``ESLH invites educators to imagine classrooms not as standardized spaces, but as soundscapes of belonging''
    \end{itemize}
    
    \item \textbf{Usage of irrelevant, loosely related, exaggerated, excessively dramatic, and poetic adjectives:}
    \begin{itemize}
        \item ``ESLH foregrounds harmony not as sameness, but as \textbf{co-resonance across difference}.''
        \item ``ESLH invites educators to imagine classrooms not as standardized spaces, but as \textbf{soundscapes of belonging}.''
        \item ``ESLH emphasizes \textbf{attunement, dignity, and healing} in response to the postcolonial, \textbf{emotional, and epistemic harms} of English.''
        \item ``\dots proposes a \textbf{praxiology} of English teaching centered on \textbf{relational presence, dialogical listening, and ethical co-creation}.''
        \item ``\dots this framework contributes to a future of language education grounded in \textbf{critical love}.''
    \end{itemize}
\end{enumerate}

The third point is more subjective, but the necessity of using terms such as ``co-resonance across difference,'' ``healing of emotional and epistemic harms of English,'' and bringing in concepts such as justice and shared humanity is highly questionable. It sounds poetic, dreamful, and even a tad apocalyptic; it lacks the impersonality and neutrality of a scientific work. It feels like the AI is doing some sort of propaganda or advertisement for the article, rather than providing a simple abstract. It sounds like the article is doing or fixing way more than it actually can.

The construction ``It is not only X, but Y'' is actually valid and really old. It is known as antithesis and was introduced by Aristotle in Rhetoric, Book III, 1410a (\url{https://kairos.technorhetoric.net/stasis/2017/honeycutt/aristotle/rhet3-9.html#1410a}). This specific form of ``negative-positive'' can even be found in the Bible, in the gospel of St. Matthew 10:34: ``Do not think that I came to send peace upon earth: I came not to send peace, but the sword'' (\url{https://www.drbo.org/chapter/47010.htm}). The difference is that in organic language, this kind of antithesis is used with a specific rhetorical motive. In the example, Jesus wanted to clarify a misconception that His words might have been causing. In the case of the article, these constructions are used very close to each other and serve little practical use. The words ``sameness'' and ``standardized spaces'' were never claims made or implied earlier in the text. No reader had assumed that harmony meant sameness, or that classrooms were, by default, standardized spaces. Such sentences happen twice within just three sentences of this abstract.

The most interesting part for me is that it was really tough for the professor to identify the moves. The abstract ended up so chaotic that the moves appeared mixed without any logic, continuity, or transition between them. This abstract was all over the place as can be seen by the highlights in \ref{app:sharedabs}.

As an extra note, I pasted the text into ZeroGPT (\url{https://www.zerogpt.com/}) to test for signs of AI-written text, and it seems to agree with my conclusion. Of course, AI detection tools are not reliable at all, but when paired with my previous analysis, it can strengthen the idea that this abstract is indeed AI-generated.

This analysis made me hypothesize that the moves may have an unintended purpose: that of identifying organic texts written by humans. If it seems that the moves and logic are scattered all around with no rhyme or reason, it can be possible evidence that the text was written by a machine imitating human behavior in abstracts without actually being able to grasp what is behind it. Exactly how AI works.

One may even say that such texts are creepy or uncanny, given how similar they are to human writing; however, when analyzed technically, there is always a sense that something is wrong. It is a text that falls under the so-called ``uncanny valley'' (a deep dive into the concept can be found here: \url{https://ieeexplore.ieee.org/abstract/document/6213238/}). This is interesting because we usually associate the uncanny valley with images and sounds that are weird, but here, such productions are texts.

We checked many other articles, and none used as many clichés as I listed in points 1, 2, and 3 above. We have many different listings with two or four items, as well as a considerable lack of sentences in the form of ``It is not only X, but Y'' and similar structures.

A general view of real-life articles showed how useful the moves really are. Even though authors often don't follow the moves too closely, they use the most important ones, and the analysis helps to understand the work and mind of the author before reading the work itself.

Before the end of the class, the professor shared that some people see children's books as a minor type of literature. She even cited an author who dismissively called children's books ``\textit{livrinhos}'', responding with: ``I'm sorry, but I prepare the readers of your books''. This is correct. Although I presented some critiques regarding the use of these books in our classes, I firmly believe they are a significant form of literature. The books we read tend to be really well made; they teach important lessons for children, help with imagination, abstract thinking, grasping fiction, and much more. They represent a valid and highly important form of art that demands a lot of studying and planning. Children think differently from adults, making it remarkably difficult to understand them well enough to write something useful for them. Furthermore, there is a huge responsibility involved; such books are part of foundational education during a time when children are being raised and introduced to the world. The language and content must be appropriate, interesting, fun, beautiful, educational, inspiring, and simple to understand. I cannot comprehend how someone would consider it easy to write stories for children. The only point I understand is that these books are primarily designed for younger audiences; for adults, philosophy and literature can provide challenges that align better with their stage in life. After all, as the author said, these books \textbf{prepare} the readers so their literature can increase in complexity.

The professor also said, ``If you do not think that the story sessions are fun, tell me so I can stop.'' Two or three of my classmates indicated that they find the stories fun. I personally find them unentertaining. Some of the artwork looks good. In general, they are repetitive. We need to stop and describe in detail what is happening on each page out loud, and the final message could easily be summarized at the beginning. I think I might be in the minority on this point. I concede that the readings are light and allow for the more talkative students to express themselves and comment on what is happening. Meanwhile, more introverted students like me can only watch and write about what is being discussed. This does not align with my definition of ``having fun''.

The most enjoyable part is reflecting on the message the book intends to communicate. To achieve that end, I would rather read the message in the form of a well-structured argument instead of going through a long \textit{reading-and-describing-what-you-see} session. Having the time to write a text reflecting on that idea during class, and then receiving feedback from the professor, would allow me to improve my argumentation and general writing while expressing my opinions on different topics. I think that scenario would fit my definition of fun much better. In this context, I mean ``fun'' as an intellectual engagement that provokes satisfaction and the possibility of practical self improvement.


\section[June 15]{June 15\textsuperscript{th}}
Today's story is called ``The Proudest Blue''; we started with a discussion on how to define pride and being proud. St. Thomas Aquinas defines pride as ``an inordinate, excessive love of one's own excellence''. Of course, the book attempts to convey a different sense. Aquinas calls this idea magnanimity: ``having noble goals, striving for excellence, and acknowledging your own worth''. It is interesting how modern society in general chose the name of the vice ``pride'' to represent this idea. I see it as a bit dangerous; it facilitates the possibility of mixing up the ideas. Perhaps using different words would be helpful to distinguish virtue from vice.

The book is useful for teaching kids about different cultures and traditions, like the Muslim culture, so if one day they see someone wearing a hijab, they will understand why they are wearing it and respect it: ``Some people won't understand your hijab, but if you understand who you are, one day they will understand too.'' Besides teaching kids about the hijab, it also serves to incentivize kids to value their own cultures and respect others.

The author's message at the end states that the book was made for people who look like her (Muslims and people of color). Honestly, I expected the idea to be more general, regarding the need to respect different cultures as a whole, moving beyond the specific focus on Muslims and people of color.

``Children can get corrupted by society.'' Then the professor discussed prejudice, noting how Black people and Haitian people suffer more than white people. Additionally, the professor brought up the point about suffering prejudice for being a nerd or for not being good at sports (being called a ``pamonha'').

She asked, ``Is it easy to be a boy in Brazil?''. That is an interesting question to receive. Although I am not the biggest fan of discussing such topics, being a man undeniably has its own challenges, pressures, prejudices, and disadvantages that tend to be ignored. The professor said, ``We don't need to be equal; we just need to learn how to deal with diversities''. This sounds correct. It is impossible for two people to be exactly alike, and some diversities demand different treatment. An example is noticing the different ideas of politeness and etiquette in Brazil, Japan, Germany, and India. What is polite in one country can be rude in another; for instance, the Japanese aversion to touching can be seen as rude by a more affectionate Brazilian.

We then continued to check memes. This meme comment session made me think about the amount of descriptions that were done in this class: looking at images, describing what is in the images, and explaining them out loud.

Then we talked about a really useful concept being introduced in the academic world called impact factor. The main idea is to give more focus to the amount of citations a researcher has. I am really happy with this new approach. It worries me that scholars in Brazil need to struggle so much to write a high quantity of publications, many times recycling the same research into different, smaller academic productions across different journals. This sounds detrimental; it feels like it can overflow the national academia with a high quantity of subpar or less useful works. Perhaps by focusing on relevance, researchers will have incentives to work on larger projects, publishing fewer works that are ultimately much more useful.


\section[June 22]{June 22\textsuperscript{nd}}
Today's book is called ``The Boy Who Grew a Forest''. It is about a boy in India who liked trees and had to deal with floods, which destroyed the vegetation. He wanted to plant more plants to help the animals and humans that suffered from the floods and the drainage. The boy grew into a man and kept planting and protecting the plants and animals from people.

This is the book with the prettiest illustrations yet.

The professor argues that this story is not far removed from Academic Writing. There are topics like reforestation, ecological restoration, environmental education, sustainability, community engagement, children's literature, environmental awareness, and the representation of real historical figures in picture books. This is different from usual, usually the professor asks ``What is the story about?''. I appreciate the analogy to academic work. But then again, doesn't almost anything relate to academic production? I can look at an apple and think about ecology, photosynthesis, nutrition and much more. It is indeed interesting to see the world in terms of academic subjects and fields of study. A good exercise to be done during class could be to pick any random thing, idea or story and write a small text following an academic style about it to understand how things in the world can be described academically and help the students to be able to explain the world in such terms.

Then we repeated the reflection about how our questions were pruned; we were not incentivized to keep asking ``why''. The idea that ``Teachers think they are hired to speak and students are to listen'' is criticized by the professor. There is a defense that teachers should let students speak and question things. That is why interruptions and curiosities are answered and accepted, instead of the teacher just saying ``I'll answer later'' and the answer never coming. This is a form of teaching people to be creative by letting them speak; she also criticizes the idea of applying the same test to everyone. It is nice to see the professor's evaluation method and teaching philosophy. I appreciate the freedom given to us students, so we can choose the articles we want to analyze, the topic we want to talk about in the paragraph, and of course, the amount of freedom we have in this Reflective Portfolio, so we can express ourselves freely in the way we better see fit. My problem would be more with the teaching philosophy, it is known that time is limited, if the professor stops to hear every opinion, comment, or answer every question, the class will simply not progress. I notice that this class has a tendency to get lost into tangents. As a student, I am significantly more interested in listening to the professor with a doctorate and an extensive scientific production and knowledge, than to talk or to listen to classmates' opinions. I never like to talk too much in class unless it is to answer a question or ask a question I see as strictly relevant and related to the topic in discussion. Otherwise it creates a sense of competition, between student and professor, and the class end up being too much oriented towards the interests of a few more talkative students, hurting the others that signed up for a class about a specific topic, a topic which only the professor has ample knowledge and authority in.

We saw some puns about punctuation. One of the puns shows a person in a trial saying ``I didn't do nothing'', the jury thinks ``oooh, a confession'', and the subtitle reads ``Jury of English Majors''. 

This pun reminds me of my studies in logic, I would like to share some of the formal logic notation to show that the joke would work equally well with a ``Jury of Logic Majors''. In propositional logic, we can formalize this sentence to demonstrate the principle of double negation. Let us define the proposition $P$ as ``I did something''. The phrase ``I did nothing'' represents the negation of $P$, formalized as $\neg P$. When the defendant says ``I didn't do nothing'', he is applying a second negation, resulting in $\neg(\neg P)$. According to the rule of double negation elimination, two negations cancel each other out, making $\neg(\neg P) \equiv P$. Logically speaking, the defendant is formally confessing to the action. 

We then talked about the serial comma (Oxford comma). The professor brought up the sentence ``Please bring Bob, a DJ and a clown'' as an argument to favor the use of the serial comma, to make it clearer that the person refers to three people. The video also brings a counterexample: ``Please bring Bob, a DJ, and a puppy''. In this example, the serial comma actually hinders understanding. People wouldn't know if they should bring one person (Bob, who happens to be a DJ) and a puppy, or two people (a DJ and Bob) and a puppy. The video talks about this distinction. Although the slides bring up this example sentence, in class the professor brought it up as another example in favor of the serial comma, with which the video does not agree.

Ironically, one person asked in class why people came up with the serial comma. The professor answered ``I don't know'' and said that this represents adult life, which is why she prefers children's life, as it is simpler.

This small comment heavily contrasts with the message from the beginning of the class about how we should behave like children and ask more ``whys''. This is what happens when we keep asking ``why'': naturally, we get deeper and deeper into a topic, so the topic gets increasingly complex, forcing us to deal with the complexity and straight-up uncertainty. Therefore, this idea of ``living like a child'' presents a heavy contradiction. If we stay curious and keep asking why, we will have to deal with answers that are increasingly complex, representing the complexity of adult life. If we choose to just accept life as it is and stop questioning, we abandon the intrinsic curiosity of childhood and embrace the conformity of adult life. Either way, living like a child is impossible. I would argue that living as an adult can be a mix of the two ways mentioned. We focus and ask the ``whys'' of a topic we wish to specialize in, choosing to stay ignorant of others because time is limited; we need to work, we get tired, and so on. 

The next topic was about comma splices, my biggest weakness. I am trying to use more dots and colons to solve that. However I always end up with either comma splices, or sounding robotic talking in topics using the dot. It saddens me a bit that at the end of the course we see these kinds of tips about commas, conjunctions, and subordinates; I felt like I could barely practice them. The same applies to the last topic of nominalization, a big characteristic of academic babble, part of how we tend to force ourselves to write in an impersonal manner. This makes the text abstract and distant. I would be happy if we saw this earlier and did some exercises in class to fix sentences and improve texts with the techniques from the TED videos.

To end the class, we watched a video arguing against the use of the words ``good'' and ``bad''. The professor adds the word ``fine'' to this list; according to her, we should be more honest in our communication. This is interesting because it provides more insight into the thought process regarding the previous activity on June 1st (when we could not repeat what the other person said). Some quotes from this discussion include: ``People are lacking real things'' and ``If you are unhappy and frustrated about being overwhelmed with tasks, well, say it!''.

I have a few problems with this perspective, starting with the implication of dishonesty when answering ``fine'' to the question ``How are you?''. We know that language serves the purpose of manifesting an intention in the intellect. We use words to externalize this intent. When someone asks me ``How are you?'', I can perfectly identify that the person does not truly intend to know exactly how I am doing. The person just wants to show courtesy and respect for me. Their intent is clear. I choose to say ``fine'' as a sign of politeness and to return the courtesy, showing that I also respect my interlocutor. Using the classical definition of a lie as the ``inadequacy of speech to the intellect'', no lie actually occurs. My intellect intends to express respect, and my speech follows with words culturally understood to achieve that exact end. Therefore honesty was completely intact in this interaction. 

The same argument can be extended to the presented video. Removing the words ``Good'' and ``Bad'' from their context in communication seems oversimplified. The examples provided in the video further demonstrate a misunderstanding of everyday social dynamics. The speaker argues that answering ``good'' when someone asks about a haircut is dishonest, in most social contexts, the person asking is simply asking for validation, not a detailed description. Just answering ``good'' or ``bad'' will be honest and exactly what they want. Even the medical example falls short. The speaker claims that telling a doctor ``I feel bad'' is a heresy that could lead to a misdiagnosis. Stating ``I feel bad'' is a generalized truth and serves as a completely valid starting point for a medical consultation. It prompts the doctor to ask follow-up questions. Ironically, the speaker suggests replacing this with ``I feel like a horde of little animals invaded my chest''. This highly subjective metaphor would be just as vague as simply saying ``bad'' and pointing to the chest, the only piece of information given in the location of the problem. The doctor would need to ask follow-up questions either way.

Building upon the classical definition of a lie as the inadequacy of speech to the intellect, one can argue that avoiding words like ``fine'', ``good'', and ``bad'' in specific social contexts is actually a form of dishonesty. When a colleague asks ``How are you?'' in passing, the intellectual intent is rarely to deliver a profound psychological evaluation. The intent is to acknowledge the social contract and express a baseline state of respect. If a person forces a highly elaborate or poetic response in this scenario, their speech no longer aligns with their simple intent. The speech becomes performative because the speaker is prioritizing a display of vocabulary over the genuine, straightforward intention of their mind. By artificially inflating the response (even if the inflation reflects a true state), the speaker creates a disconnect between the intellect and the spoken word. Therefore, demanding complex descriptions for simple intellectual states forces individuals into performative dishonesty. This strongly suggests that words like ``fine'', ``good'', and ``bad'' are the most honest and accurate choices for everyday social interactions.

From this conclusion, I would argue that demanding more complete answers from people during casual social conversations constitutes a form of dishonesty. By doing so, one deliberately chooses to ignore a perfectly valid answer in favor of performance and dissimulation. In academic writing, being specific and descriptive is required, making the avoidance of vague words necessary. That's why I understand the application of the video in class for that specific purpose. However it is important to clarify to the class that the video acts as a significant hyperbole; in day-to-day life, it remains completely honest and perfectly acceptable to use ``fine'', ``good'', and ``bad''.

\section[June 29]{June 29\textsuperscript{th}}
Today is my last day at FURB. The class bought some special ``Festa Junina'' food. I am feeling especially nostalgic. Before class, I walked through the campus, visited every building, took some pictures of the places where I studied, and truly reminisced on the past four years I spent there. Now I do not have any more classes to attend. There was no reading today; we watched a 30-minute video on the story ``The Boy, the Mole, the Fox and the Horse''. The animation and art style were really beautiful. The story is about a boy lost in the snow who makes friends with a mole, a fox, and a horse. Together, they try to find a home for the boy by following the river and then the lights. 

The movie's main characteristic lies in the sentences the characters say. They are poetic and hold deeper meanings. I will reflect on a few of them. The kid complains that even though they are walking toward the lights, they always seem further and further away. This reminds me of my own situation. I really wanted to graduate. During my path, I completed a five-month exchange program in Germany, and now I will spend two whole years studying there without formally getting my diploma. The sensation of having the finish line move further away can be quite frustrating. I know this extension is for the best, and I am ready to persevere.

It strikes me as weird the way the characters are so quick to say that they love each other. They barely know each other. This sense of instant commitment appears a bit shallow. ``Love'' is a strong word, and I believe they show a good example of friendship. St. Thomas Aquinas defines love as ``to will the good of the other'' (more information can be found here: \url{https://www.newadvent.org/summa/2026.htm#article2}). The stronger the love, the more unconditional that will becomes. This accurately represents the idea that the animals accepted the boy's desire to go home, willing the good for him despite their own sadness. The horse protecting them from the storm and the fox saving the mole are also excellent examples of this definition. The movie represents love in a truly beautiful way.

Another quote that I like from the movie is ``One of the great freedoms is how we react to things''. This conveys the notion that we cannot control everything. Many unpleasant things will happen to us. We will always remain in control of our reactions. ``Help is the bravest thing you've ever said'' is another excellent quote. Asking for help is a highly important practice to cultivate, especially for children. It requires humility and the understanding that sometimes we cannot do everything on our own, alongside the courage to admit that to another person. This practice can save lives in case someone is in danger. This is also one of the first lessons we learn when starting our careers as programmers. When we do not know something, asking for help is the best course of action to avoid causing problems for the team and the company.

As a last note, I liked the music that plays during the movie. It is calming, the script is charming, and the pictures are pretty. It is highly beneficial for children to have access to calmer cartoons, completely removing the chaotic stimuli present in other media.

One point about the Reflective Portfolio. I wrote it using a really popular tool in academic writing in exact and natural sciences called \LaTeX. Using it I could ``code'' my paper. If there is curiosity, the source code for the paper can be found here: \url{https://github.com/GuerreiroG/reflective-portfolio-aw}

Afterward, the class followed the orientations for the conclusion, so I will address those questions in the final section.